\documentclass[11pt]{article}
\usepackage[margin=1in]{geometry}
\usepackage{amsmath, amssymb, amsthm}
\usepackage{booktabs}
\usepackage{enumitem}
\usepackage{hyperref}

\newcommand{\R}{\mathbb{R}}
\newcommand{\E}{\mathbb{E}}
\newcommand{\Var}{\mathrm{Var}}
\newcommand{\Unif}{\mathrm{Unif}}
\newcommand{\Cov}{\mathrm{Cov}}
\newcommand{\W}{W}
\newcommand{\Norm}{\mathcal{N}}
\newcommand{\Fate}{F_{\mathrm{ATE}}}
\newcommand{\Flin}{F_{\mathrm{lin}}}

\title{Optimal-Transport Aggregation of CATE Distributions \\ into an ATE Distribution}
\date{}

\begin{document}
\maketitle

\section{Problem statement}

We have $N$ units. For each unit $i \in \{1, \dots, N\}$ a foundational prior network has produced a 2D joint density of the potential outcomes $(Y_i(0), Y_i(1))$, which we have fit with MALC\_2D:
\[
  p_i(y_0, y_1) \quad \text{on the same 2D grid for all } i.
\]
The conditional average treatment effect (CATE) for unit $i$ is the random variable
\[
  D_i \;=\; Y_i(1) \;-\; Y_i(0), \qquad D_i \sim f_i(d).
\]
The goal is to aggregate $\{f_1, \dots, f_N\}$ into a single \textbf{population ATE distribution} $\Fate(d)$ that respects the \textbf{modal / causal-world structure} of the inputs: when each $f_i$ is a mixture over $K$ latent ``causal worlds'', $\Fate$ should also have $K$ modes, with each mode being the within-world average of the corresponding mode across units.

\section{Notation}

\begin{center}
\begin{tabular}{ll}
\toprule
Symbol & Meaning \\
\midrule
$i = 1, \dots, N$ & unit (patient) index \\
$k = 1, \dots, K$ & causal-world (mode) index \\
$p_i(y_0, y_1)$ & 2D joint density from MALC\_2D \\
$f_i(d)$ & 1D CATE density for unit $i$ \\
$\pi_{i,k}$ & mixing weight of world $k$ for unit $i$, $\sum_k \pi_{i,k} = 1$ \\
$g_{i,k}(d)$ & within-world density of unit $i$ for world $k$ (log-concave) \\
$G_{i,k}(d)$ & CDF of $g_{i,k}$ \\
$Q_{i,k}(\tau) = G_{i,k}^{-1}(\tau)$ & quantile function of $g_{i,k}$ \\
$g_k^\star(d)$ & per-world 1D 2-Wasserstein barycenter \\
$w_k$ & aggregate weight of world $k$, $\sum_k w_k = 1$ \\
$\Fate(d)$ & final ATE distribution \\
$\W_2$ & 2-Wasserstein distance \\
\bottomrule
\end{tabular}
\end{center}

\section{Deriving the 1D CATE density from the 2D joint}

For each unit $i$, change variables from $(Y_0, Y_1)$ to $(Y_0, D)$ with $D = Y_1 - Y_0$. The Jacobian determinant is $1$, so the joint density of $(Y_0, D)$ is
\[
  \tilde p_i(y_0, d) \;=\; p_i(y_0, y_0 + d).
\]
Marginalising out $y_0$ gives the CATE density
\begin{equation}
  \boxed{\;f_i(d) \;=\; \int_{\R} p_i\!\bigl(y_0,\; y_0 + d\bigr)\, dy_0\;}
  \label{eq:cate-density}
\end{equation}
Numerically, we sample $y_0$ on a fine 1D grid covering the range of $p_i$ and approximate the integral by a Riemann sum along anti-diagonal lines $y_1 = y_0 + d$.

\section{Per-unit decomposition with MALC\_BM}

For a chosen $K$, fit each unit's CATE density with MALC\_BM:
\begin{equation}
  f_i(d) \;=\; \sum_{k=1}^{K} \pi_{i,k}\, g_{i,k}(d),
  \qquad \pi_{i,k}\ge 0, \quad \sum_{k=1}^{K}\pi_{i,k} = 1,
  \label{eq:decomp}
\end{equation}
where each $g_{i,k}$ is a \textbf{log-concave 1D density} (the MALC\_BM per-component fit). After fitting, sort components within each unit by mode location so that the index $k$ refers to the same world across all $i$:
\[
  \operatorname{mode}(g_{i,1}) < \operatorname{mode}(g_{i,2}) < \cdots < \operatorname{mode}(g_{i,K}).
\]

\section{1D 2-Wasserstein barycenter --- closed form}

Given $M$ probability measures on $\R$ with densities $\nu_1, \dots, \nu_M$ and weights $\alpha_1, \dots, \alpha_M$ ($\alpha_m \ge 0$, $\sum_m \alpha_m = 1$), the \textbf{2-Wasserstein barycenter} is
\[
  \nu^\star \;=\; \arg\min_{\mu} \;\sum_{m=1}^{M} \alpha_m\, \W_2^2(\mu,\, \nu_m).
\]
In one dimension this has a closed form via quantile functions. If $Q_m(\tau) = F_m^{-1}(\tau)$ is the quantile function of $\nu_m$, then the barycenter has quantile function
\begin{equation}
  \boxed{\;Q^\star(\tau) \;=\; \sum_{m=1}^{M} \alpha_m\, Q_m(\tau), \qquad \tau \in (0, 1).\;}
  \label{eq:bary-1d}
\end{equation}
The barycenter density $\nu^\star$ is recovered from $Q^\star$ by inversion: $\nu^\star$ is the pushforward of the uniform measure on $(0,1)$ by $Q^\star$.

This is a special case of Agueh \& Carlier (2011) and has the property that if all $\nu_m$ are translations of a common shape $\nu_0$ --- $\nu_m(x) = \nu_0(x - c_m)$ --- then $\nu^\star(x) = \nu_0\bigl(x - \sum_m \alpha_m c_m\bigr)$, i.e.\ the barycenter is the canonical shape at the mean translation. Linear (mixture) averaging would \emph{blur} the shape across translations.

\section{Per-world barycenter}

For each world $k$, define the \textbf{renormalised within-world weights}
\[
  \alpha_{i,k} \;=\; \frac{\pi_{i,k}}{\sum_{j=1}^{N} \pi_{j,k}}, \qquad i = 1, \dots, N,
\]
so $\sum_i \alpha_{i,k} = 1$. The per-world 1D Wasserstein barycenter is
\begin{equation}
  \boxed{\;Q_k^\star(\tau) \;=\; \sum_{i=1}^{N} \alpha_{i,k}\, Q_{i,k}(\tau),
  \qquad g_k^\star(d) \;=\; \bigl(Q_k^\star\bigr)_\# \Unif(0,1).\;}
  \label{eq:perworld-bary}
\end{equation}
In words: at every probability level $\tau$, the world-$k$ barycenter quantile is the $\pi$-weighted average of the world-$k$ quantiles of each unit.

\section{Aggregate ATE distribution}

The aggregate weight of world $k$ across the sample is
\[
  w_k \;=\; \frac{1}{N}\sum_{i=1}^{N} \pi_{i,k}, \qquad \sum_{k=1}^{K} w_k = 1.
\]
The final ATE distribution is
\begin{equation}
  \boxed{\;\Fate(d) \;=\; \sum_{k=1}^{K} w_k\, g_k^\star(d).\;}
  \label{eq:fate}
\end{equation}
This is a $K$-mode density: world $k$ contributes a single mode, located at the within-world average mode, with mass $w_k$ equal to the empirical mean of the per-unit world-$k$ weights.

\section{Comparison with the naive linear mixture}

The pointwise (Monte-Carlo / linear-mixture) aggregator is
\[
  \Flin(d) \;=\; \frac{1}{N}\sum_{i=1}^{N} f_i(d).
\]
Substituting the MALC\_BM decomposition of each $f_i$,
\[
  \Flin(d) \;=\; \frac{1}{N}\sum_{i=1}^{N}\sum_{k=1}^{K} \pi_{i,k}\, g_{i,k}(d)
              \;=\; \sum_{k=1}^{K} w_k\, h_k(d),
\]
where the within-world \textbf{linear mixture} is
\[
  h_k(d) \;=\; \sum_{i=1}^{N} \alpha_{i,k}\, g_{i,k}(d).
\]
The aggregate world weights $w_k$ are identical; the difference is purely within-world:
\begin{equation}
  \Fate(d) - \Flin(d) \;=\; \sum_{k=1}^{K} w_k \bigl[\, g_k^\star(d) - h_k(d) \,\bigr].
  \label{eq:diff}
\end{equation}
If all $g_{i,k}$ coincide ($g_{i,k} = g_k$ for every $i$), then $g_k^\star = h_k = g_k$ and the two aggregators agree. Otherwise --- when there is within-world heterogeneity in shape or location --- $h_k$ is a Wasserstein-smeared version of $g_k^\star$ and is wider than $g_k^\star$ (variance is additive in the linear mixture, while the Wasserstein barycenter averages quantiles, which is variance-preserving under pure translation and shrinks within-world variance under shape mismatch).

\section{Why we cannot skip the decomposition step}

A natural question: can we just compute the 1D Wasserstein barycenter of the $\{f_i\}$ directly, without the MALC\_BM decomposition?

\emph{No}, because the quantile-averaging barycenter does not preserve mixing weights when those weights differ across inputs. Consider two bimodal CATE densities with the \emph{same} mode locations but different weights:
\[
  f_1(d) = 0.5\,\Norm(-2, \sigma^2) + 0.5\,\Norm(+2, \sigma^2), \qquad
  f_2(d) = 0.3\,\Norm(-2, \sigma^2) + 0.7\,\Norm(+2, \sigma^2).
\]
The ``correct'' average is $F_{\mathrm{target}} = 0.4\,\Norm(-2, \sigma^2) + 0.6\,\Norm(+2, \sigma^2)$. But the quantile-averaged barycenter $Q^\star(\tau) = \tfrac12 Q_1(\tau) + \tfrac12 Q_2(\tau)$ behaves as

\begin{center}
\begin{tabular}{cccll}
\toprule
$\tau$ range & $Q_1$ & $Q_2$ & $Q^\star$ & comment \\
\midrule
$(0, 0.3)$ & $\approx -2$ & $\approx -2$ & $\approx -2$ & left mode \\
$(0.3, 0.5)$ & $\approx -2$ & $\approx +2$ & $\approx 0$ & phantom mass at 0 \\
$(0.5, 1)$ & $\approx +2$ & $\approx +2$ & $\approx +2$ & right mode \\
\bottomrule
\end{tabular}
\end{center}

The interval $\tau \in (0.3, 0.5)$ where one quantile axis ``switches modes'' before the other contributes a spurious central mode with mass $0.2$. The decomposition step circumvents this by aligning the world axis explicitly: world-1 mass is averaged with world-1 mass, world-2 mass with world-2 mass, independently of any quantile mismatch.

\section{Algorithm summary}

\begin{enumerate}[leftmargin=*]
\item \textbf{Inputs.} $N$ units each with 2D MALC\_2D fit $p_i$, a common 1D grid for $d$, the number of causal worlds $K$.

\item \textbf{Per unit $i$:}
\begin{enumerate}
  \item Compute CATE density $f_i(d) = \int p_i(y_0, y_0 + d)\, dy_0$ via anti-diagonal integration on the 2D grid.
  \item Decompose with MALC\_BM at fixed $K$: $f_i = \sum_k \pi_{i,k} g_{i,k}$. Sort components by mode location so that ``world $k$'' is consistent across units.
\end{enumerate}

\item \textbf{Per world $k = 1, \dots, K$:}
\begin{enumerate}
  \item Renormalise weights $\alpha_{i,k} = \pi_{i,k} / \sum_j \pi_{j,k}$.
  \item Compute the 1D Wasserstein barycenter quantile function $Q_k^\star(\tau) = \sum_i \alpha_{i,k} Q_{i,k}(\tau)$ on a uniform $\tau$-grid.
  \item Recover the barycenter density $g_k^\star$ from $Q_k^\star$ by inversion / interpolation.
\end{enumerate}

\item \textbf{Combine.} Aggregate world weights $w_k = \frac{1}{N}\sum_i \pi_{i,k}$ and report
\[
  \Fate(d) = \sum_{k=1}^{K} w_k\, g_k^\star(d).
\]

\item \textbf{Validation.}
\begin{itemize}
  \item Sample-size ablation: repeat for $N' < N$, compare $\Fate^{(N')}$ to $\Fate^{(N)}$ or to a known truth.
  \item Ground-truth simulation: known DGP for $(\pi_{i,k}, g_{i,k})$; report $\W_2(\Fate, F_{\mathrm{true}})$ and $\W_2(\Flin, F_{\mathrm{true}})$.
\end{itemize}
\end{enumerate}

\section*{References}
\begin{itemize}
\item Agueh, M., Carlier, G. (2011). \emph{Barycenters in the Wasserstein space.} SIAM J.\ Math.\ Anal.
\item Balabdaoui, F., Jankowski, H., Rufibach, K., Pavlides, M. (2013). \emph{Asymptotics of the discrete log-concave maximum likelihood estimator and related applications.} JRSS B. (used for \texttt{logConDiscrMLE})
\item Cule, M., Samworth, R., Stewart, M. (2010). \emph{Maximum likelihood estimation of a multidimensional log-concave density.} JRSS B. (used for the 2D MLE in MALC\_2D)
\item D\"umbgen, L., Rufibach, K. (2009). \emph{Maximum likelihood estimation of a log-concave density: basic properties and uniform consistency.} Bernoulli. (used for the 1D continuous MLE)
\end{itemize}

\end{document}
